\chapter{Overview and Scope}\label{ch:overview}

\section{Problem statement}

Month-end Trial Balance review at SCB is currently executed as a
spreadsheet-driven process. For each of the 234 Legal Entities (LEs)
in the Group, a GFS Analyst rolls the prior month working file
forward, downloads the PSGL master account extract, runs variance
calculations in an EUC-governed Excel workbook, chases country
stakeholders for commentary, tracks unposted journals, and hands a
signed-off summary to the Head of Africa GFS for inclusion in the
monthly FORTM pack. The cycle runs from \texttt{M-3} (preparation) to
\texttt{M+7} (cycle close), with the critical review phase compressed
into \texttt{M-2} through \texttt{M+5}.

Three problems follow from the current model:

\begin{enumerate}
  \item \textbf{Effort.} The process consumes 20--25 FTEs at roughly
        four hours/day during the active review window --- an
        establishment the business case identifies as the largest
        single automation candidate in the close calendar.
  \item \textbf{Control surface.} Three EUC references govern the
        process (\path{EX/200/1353107/0001/000}, variance file
        \path{EX/200/1579702/0001/000}, and review control
        \path{EX/200/1579702/0002/000}). EUCs carry recurring
        evidencing burden and are misaligned with the UK ACG ambition
        to eliminate end-user spreadsheets.
  \item \textbf{Latency.} The monthly batch cadence rules out daily
        monitoring; anomalies surface at month-end rather than when
        they arise.
\end{enumerate}

\section{In-scope workstreams}

This specification covers the three workstreams confirmed with the
business sponsor:

\begin{description}
  \item[A. Canonical corpus.] A machine-readable representation of
        every artefact under review --- the 28-state BPMN workflow,
        five decision points, three EUC controls, six KPIs, eight
        process steps, and the domain glossary --- normalised to JSON
        and validated against Draft-07 schemas.
  \item[B. AI augmentation.] An on-premise pipeline hosted on the
        \texttt{vllm-only} policy (\texttt{data-product.routing = vllm-only})
        that ingests the PSGL / S/4 extract, applies
        \emph{statistical} anomaly detection (percentile-based
        thresholds, configurable per account category) at caption,
        segment, and LE levels, drafts commentary using historical
        patterns, and emits the FORTM narrative. Machine-learning
        anomaly detectors (Isolation Forest, autoencoder, Prophet)
        are scoped for v2.0 (Chapter~\ref{ch:future-ml}) and are
        not part of the v1.0 deliverable.
  \item[C. Workflow and approval engine.] A formalisation of the seven
        DOI activities and the gateway logic that decides when
        material variances require commentary, when journals must be
        posted, when further investigation is needed, and whether the
        entry-confirmation deadline at \texttt{M+4} has been honoured.
\end{description}

\section{Out-of-scope}

The following are referenced but not specified:

\begin{itemize}
  \item Country IFRS reporting --- handled by local Finance teams.
  \item BSS Risk Assessment --- sibling process with its own business
        case. Its source document is retained in the corpus for
        terminology reuse only.
  \item SAP / PeopleSoft integration detail --- the AI layer reads
        data through the same \texttt{SC Bridge} interface as the
        manual process; interface engineering is a separate document.
  \item Country expansion beyond the Hong Kong (HK) pilot --- the
        HKG\_TB and HKG\_PL workbooks are the source of truth for the
        initial rollout; LE coverage extends to all 234 entities in
        phase~4 onward.
\end{itemize}

\section{Stakeholders}

\Cref{tab:stakeholders} enumerates the stakeholders named in the
business case and the DOI.

\begin{table}[htbp]
\caption{Stakeholders and responsibilities.}\label{tab:stakeholders}
\centering
\small
\begin{tabularx}{\linewidth}{@{}llX@{}}
\toprule
\textbf{Role} & \textbf{Named} & \textbf{Responsibilities} \\ \midrule
Process Owner             & Modi, Piyush & BPMN ownership; workflow changes \\
Head of Africa GFS        & ---          & Internal review; FORTM sign-off \\
GFS Analyst               & ---          & Primary executor; variance file owner \\
Country Stakeholders      & group        & Commentaries; \texttt{M+4} entry confirmation \\
GCFO Teams                & ---          & Consume narrative; decision-making \\
Group Team                & ---          & FORTM recipient; further review \\
AI Architect              & ---          & PSGL + S/4 compatibility \\
PPE / AI Architect        & ---          & SAP integration assessment \\
\bottomrule
\end{tabularx}
\end{table}

\section{Timeline}

The process obeys a fixed month-end-relative timeline
(\cref{tab:timeline}). All state transitions in \cref{ch:workflow}
carry a \texttt{timing} field that maps to one of these phases.

\begin{table}[htbp]
\caption{Month-end-relative timeline phases.}\label{tab:timeline}
\centering
\small
\begin{tabularx}{\linewidth}{@{}llX@{}}
\toprule
\textbf{Phase} & \textbf{Timing} & \textbf{Activities} \\ \midrule
Preparation & M-3        & Roll forward; PSGL download; master refresh; parameter set \\
Extraction  & M-3 to M+5 & SC Bridge extract into the variance file (\texttt{EUC-002}) \\
Analysis    & M-2 to M+5 & Variance calc; commentary; investigation; risk \\
Review      & M+5        & Internal review (\texttt{EUC-003}); FORTM submission \\
\bottomrule
\end{tabularx}
\end{table}

\section{Context diagram}

\begin{figure}[htbp]
\centering
\begin{tikzpicture}[
  node distance=8mm and 10mm,
  box/.style={draw,rounded corners=2pt,minimum height=9mm,minimum width=26mm,
              align=center,font=\small},
  sys/.style={box,fill=tbblue!15,draw=tbblue!60},
  ppl/.style={box,fill=tbgreen!15,draw=tbgreen!60},
  outp/.style={box,fill=tbamber!15,draw=tbamber!60},
  arr/.style={-{Latex[length=2mm]},thick}
]
  \node[sys] (psgl)   {PSGL / S/4};
  \node[sys,right=of psgl] (bridge) {SC Bridge};
  \node[sys,right=of bridge] (xl)  {Variance file \\ \scriptsize \texttt{EUC-002}};
  \node[ppl,below=of psgl] (gfs)   {GFS Analyst};
  \node[ppl,below=of bridge] (cty) {Country Stakeholders};
  \node[ppl,below=of xl]    (hd)   {Head of Africa GFS};
  \node[outp,right=18mm of hd] (fortm) {FORTM pack \\ \scriptsize Group};

  \draw[arr] (psgl)   -- (bridge);
  \draw[arr] (bridge) -- (xl);
  \draw[arr] (gfs)    -- (xl);
  \draw[arr] (cty)    -- (xl) node[midway,right,font=\scriptsize]{commentary};
  \draw[arr] (xl)     -- (hd);
  \draw[arr] (hd)     -- (fortm);
\end{tikzpicture}
\caption{Context of the TB review process: three systems (blue), three
human roles (green), and one downstream output (amber). The
variance-analysis workbook is the primary EUC artefact.}\label{fig:context}
\end{figure}

\section{Expected benefits}

\begin{itemize}
  \item FTE savings: 5 FTEs redeployed (business-case field
        \texttt{executiveSummary.\allowbreak fteImpact.\allowbreak expectedSavings}).
  \item Retirement of three EUCs (Group ACG alignment) on a
        phased schedule: the three Excel-based EUCs
        (\texttt{EUC-001}, \texttt{EUC-002}, \texttt{EUC-003}) are
        replaced by the rule-firing log plus the review dashboard
        audit trail (Chapter~\ref{ch:review-dashboard}). Shared
        Drive evidence storage is retained only as a read-only
        archive for a one-year deprecation window
        (see Section~\ref{sec:euc-retirement-plan}); the dashboard
        audit trail is the authoritative evidence source from
        v1.0 cutover onwards.
  \item Enablement of daily controls (from monthly batch).
  \item Foundation for the SCB Finance LLM.
\end{itemize}

\section{EUC retirement plan}\label{sec:euc-retirement-plan}

The three EUCs are retired, not merely supplemented, on the
following schedule. ``Retirement'' means the EUC is no longer the
authoritative evidence for ACG sampling; the dashboard audit trail
(Chapter~\ref{ch:review-dashboard}, \cref{sec:dashboard-audit-trail}) is authoritative.

\begin{table}[htbp]
\centering
\caption{EUC retirement schedule.}
\label{tab:euc-retirement}
\footnotesize
\begin{tabularx}{\linewidth}{@{}l l l X@{}}
\toprule
\textbf{EUC} & \textbf{Artefact} & \textbf{Retirement cutover} & \textbf{Successor evidence} \\ \midrule
EUC-001 & PSGL master extract & Month of v1.0 go-live & Dashboard ``Master Data'' tab; SHA-256 of extract in audit trail \\
EUC-002 & Variance analysis workbook & Month of v1.0 go-live & Dashboard ``Variance Review'' tab; per-item decision record \\
EUC-003 & FORTM Pack & Month of v1.0 go-live + 1 cycle & Dashboard ``FORTM Export'' with signed PDF snapshot \\
\bottomrule
\end{tabularx}
\end{table}

\noindent
During the deprecation window (12 months from v1.0 go-live), the
Shared Drive evidence folder is mirrored automatically from the
dashboard export at month-end (read-only, with mismatch alerting).
After the deprecation window closes, the Shared Drive location is
decommissioned and removed from the ACG sampling procedure.
Section~\ref{sec:controls-eucs} in Chapter~6 is superseded by this
plan for all v1.0 and later cycles; the Chapter~6 text is retained
as a historical reference only.
